\section{Related Work}

\paraf{Optimization for RL framework.}
The benefits of LLM RL have drawn great attention to RL framework optimizations~\citep{yao2023deepspeed,sheng2025hybridflow,hu2024openrlhf,zhong2025optimizing,zhong2025streamrl,fu2025areal,he2025history,wang2025distflow,lei2024puzzle,shen2024nemo,han2025asyncflow}.
OpenRLHF~\citep{hu2024openrlhf} and verl~\citep{sheng2025hybridflow} explore the disaggregated and colocated architectures separately.
Some following frameworks propose fusion~\citep{zhong2025optimizing}, pipelining~\citep{zhong2025optimizing}, streaming~\citep{zhong2025streamrl}, and asynchronous execution~\citep{fu2025areal} to improve system efficiency.
Besides, some recent optimizations propose request-level scheduling, balancing, and speculative decoding~\citep{gao2025rollpacker,jinefficient},  to further speed up the rollout phase.
Despite these advancements, existing frameworks often treat external invocations as secondary, either overlooking their resource management or defaulting to static over-provisioning. 
This leads to significant resource waste, particularly because agentic RL tasks involve longer trajectories and more complex environment interactions.

\parabf{Agent resource management.}
The increasing complexity of agentic RL has necessitated more sophisticated management of external resources used for tools and rewards.
Recent efforts like Mars~\citep{zhu2026mars} have addressed reward computation efficiency by deploying it as a standalone service.
Mars leverages a "request-in, batch-out" characteristic to relax request-level latency requirements into batch-level constraints.
While effective for rewards, \sysname provides a more comprehensive solution by establishing a unified action-level platform for all types of tool usage and reward services, including API calls, code execution, and simulators.
% By shifting control from long-lived trajectories to individual atomic interactions (actions), \sysname enables statistical multiplexing across different tasks and reduces resource idle time that a single service may miss.
MegaFlow~\citep{zhang2026megaflow}, SkyRL-Agent~\citep{cao2025skyrl}, and ROCK~\citep{wang2025let} abstract agent service and adopt techniques in cloud computing to manage tools.
By integrating tool characteristics like scalability and sharing pattern, \sysname enables specialized optimization for tool invocations and reward services.

\parabf{Resource auto-scaling.}
Autoscaling is a cornerstone of cloud computing, employing reactive, proactive, or hybrid approaches~\citep{lorido2014review,wu2023xron,wang2024autothrottle,wu2024can,zhang2024jolteon} to maintain service-level objectives (SLOs).
In the LLM domain, systems~\citep{zhang2025blitzscale,yu2025lambda,xiang2025aegaeon,zhao2025seallm} have optimized autoscaling for large-scale inference workloads.
For RL-specific tasks, Mars~\citep{zhu2026mars} introduces history-based scaling and timeout-aware mechanisms to predict and adjust resource requirements at batch boundaries.
\sysname extends these principles by introducing an elastic action-level scheduler designed to handle the bursty and heterogeneous nature of agentic RL actions.
By modeling the scalability of various actions, \sysname minimizes the total ACT and thus, accelerates the RL training process.
